% Member 1 report sections: inventory and stock integrity (HLD 5 service 3, HLD 7 data design,
% HLD 9 scalability and consistency, one HLD 13 trade-off).
% Body only, so the team report can % Member 1 report sections: inventory and stock integrity (HLD 5 service 3, HLD 7 data design,
% HLD 9 scalability and consistency, one HLD 13 trade-off).
% Body only, so the team report can % Member 1 report sections: inventory and stock integrity (HLD 5 service 3, HLD 7 data design,
% HLD 9 scalability and consistency, one HLD 13 trade-off).
% Body only, so the team report can % Member 1 report sections: inventory and stock integrity (HLD 5 service 3, HLD 7 data design,
% HLD 9 scalability and consistency, one HLD 13 trade-off).
% Body only, so the team report can \input{member1-inventory}. Needs: array, booktabs.
% member1-standalone.tex compiles these sections on their own.
% Numbers measured on branch m1/postgres-audit-load-test (on top of main 7c95078), 21 November 2026, one laptop.

\section{Inventory (HLD 5, service 3)}
\label{sec:m1-inventory}

Every other service trusts one number from Inventory: how many units of a product a store can still sell. If that
number is too high, a customer pays for something that is not there. Inventory therefore answers two questions and
guards one rule.

\begin{itemize}
  \item \textbf{How many can we sell?} Available to promise (ATP) is on hand, minus held, minus a safety buffer, and
    never below zero. The buffer hides a few units of a product from customers so that a small counting error does
    not become a failed order. Replenishment and picking ignore it, because those units are really on the shelf.
  \item \textbf{Where are they?} Each batch of a product in each bin is one row with its expiry date. Picking takes the
    batch that expires first, then the nearest bin, and never a batch that has already expired.
  \item \textbf{The rule: a unit is promised to one customer only.} Checkout \emph{holds} units, packing
    \emph{commits} the sale, and everything in between only has to give a hold back.
\end{itemize}

\paragraph{Hold, then commit.} Checkout raises \texttt{qty\_reserved} and leaves \texttt{qty\_on\_hand} alone. Only
packing lowers both. A failed payment, a cancelled order or an expired hold therefore never has to ``put stock
back''; it lowers \texttt{qty\_reserved} and nothing else. A hold lasts 10 minutes. A timed job releases expired
holds every 30 seconds, and a hold is released by one conditional \texttt{UPDATE} on its status, so a cancel, a
payment and the timed job arriving together cannot give the same units back twice.

\paragraph{What else the service does.}
\begin{itemize}
  \item \textbf{Expiry.} An hourly job writes off every batch past its date, with one ledger row per batch, and
    publishes \texttt{item.near\_expiry} for batches within two days of their date, once per batch per day.
  \item \textbf{Idempotency keys.} A stock adjustment sent with an \texttt{Idempotency-Key} header changes stock
    once; a retry gets the first answer back. The key is saved in the same transaction as the stock change, so a
    crash between the two cannot separate them.
  \item \textbf{Cycle counts.} A counter is told which bin to count but not what the records expect. A difference
    worth at most Rs~500 is applied at once; a larger one waits for a manager, who may not approve their own count.
    When a picker reports an empty shelf, a count of that bin is opened automatically. The share of counted lines that
    matched is the inventory accuracy figure of HLD~12.
  \item \textbf{Functions for the other services.} Fulfilment asks Inventory to choose bins, to commit a sale and to
    write off units a picker could not find. Inbound asks it to receive stock. Two more were added for work the
    other members have planned: \texttt{swap\_line} holds a substitute for a dropped line, and \texttt{move}
    shifts units between two bins of the same temperature zone. No other module writes to the stock tables.
\end{itemize}

\section{Data design (HLD 7)}
\label{sec:m1-data}

\paragraph{Three records that must agree.} Stock is written in three places, and
Table~\ref{tab:m1-stock-records} says what each is for. For every store and product, the total in
\texttt{stock\_levels}, the sum of the bins in \texttt{stock\_items} and the sum of the movements in
\texttt{stock\_ledger} must be equal. Every function that changes one changes all three in the same transaction.

\begin{table}[ht]
  \centering
  \footnotesize
  \begin{tabular}{@{}>{\raggedright\arraybackslash}p{0.17\linewidth}>{\raggedright\arraybackslash}p{0.25\linewidth}
                  >{\raggedright\arraybackslash}p{0.5\linewidth}@{}}
    \toprule
    Table & One row per & Used for \\
    \midrule
    \texttt{stock\_levels} & store and product & On hand, held and safety buffer. Every checkout hold is one
      conditional \texttt{UPDATE} on this row \\
    \addlinespace[3pt]
    \texttt{stock\_items} & store, product, bin and batch & Where the units are and when they expire: pick paths,
      first-expiry-first-out, putaway, cycle counts \\
    \addlinespace[3pt]
    \texttt{stock\_ledger} & stock movement & Append-only history with a reason (receipt, sale, damage, expiry,
      theft, count, not found, move). A mistake is undone by a new row with the opposite sign, never by an edit \\
    \bottomrule
  \end{tabular}
  \caption{The three records of stock.}
  \label{tab:m1-stock-records}
\end{table}

HLD~7.1 keeps the held quantity on every batch row. The prototype keeps it on a separate totals row instead. A
checkout does not care which bin a unit is in, so the hold should lock one small row per product, not several batch
rows in an order that could deadlock. The cost is that the totals row repeats what the batch rows already say, and
two copies can drift apart.

\paragraph{The audit.} \texttt{GET /v1/admin/stock-audit} compares the three records for every store and product
and lists each disagreement with its values: level against bins, level against ledger, a negative quantity, or
more units held than exist. An empty list means the stock is consistent. The automated tests call the same function
after every test that moves stock, so a change that breaks the rule fails in review, not in a store.

\paragraph{PostgreSQL and migrations.} HLD~7.2 chooses PostgreSQL because its transactions and row locks prevent
overselling. The prototype now runs on either database, chosen by one setting, \texttt{DARKSTORE\_DB}: a SQLite file
for a laptop with nothing to install, PostgreSQL for a hosted deployment. All 313 automated tests pass on both
(SQLite and PostgreSQL~16), and the build server runs them on both for every change. Table changes are Alembic
migrations: a change to the models is turned into a migration file with one command and applied with a second, and
the app applies pending migrations itself when it starts, so a hosted database is changed in place and keeps its
data. A test fails when the models and the migrations disagree.

\section{Scalability and consistency (HLD 9)}
\label{sec:m1-consistency}

\subsection{No overselling: check and change in one statement}

The hold is a single SQL statement whose \texttt{WHERE} clause is the stock check:
\begin{center}\small
  \verb|UPDATE stock_levels SET qty_reserved = qty_reserved + :qty|\\
  \verb|WHERE store_id = :store AND sku_id = :sku|\\
  \verb|  AND qty_on_hand - qty_reserved - safety_buffer >= :qty|
\end{center}
If the statement changes no row, there was not enough stock and the whole cart is rolled back, so a cart is held
whole or not at all. Reading the level first and updating it afterwards would be wrong: two requests could both read
``1 left'' and both take it.

The statement is safe on PostgreSQL at its default isolation level (read committed) for a specific reason. When two
requests update the same row, the second waits for the first to commit, and then \emph{checks its \texttt{WHERE}
clause again against the new row} before changing anything. The second buyer of the last unit therefore sees zero
available and changes no row. SQLite reaches the same result more bluntly: it allows one writer at a time.

\subsection{The proof: 200 buyers, 50 units}

\texttt{scripts/load\_checkout.py} resets the data, puts exactly 50 units of one product on the shelf, and lets 200
customers each try to hold one unit at the same instant against a running server. It passes only if exactly 50 holds
succeed, the other 150 buyers are told \texttt{409 insufficient\_stock}, and the audit afterwards is empty.
Table~\ref{tab:m1-race} gives the result of five runs on each database.

\begin{table}[ht]
  \centering
  \footnotesize
  \begin{tabular}{@{}lrrrrrr@{}}
    \toprule
    Database & Holds & Refused & Errors & Audit mismatches & p50 (ms) & p99 (ms) \\
    \midrule
    SQLite (one file, WAL mode) & 50 & 150 & 0 & 0 & 883 & 1998 \\
    PostgreSQL 16 & 50 & 150 & 0 & 0 & 1337 & 1478 \\
    \bottomrule
  \end{tabular}
  \caption{Checkout race: 200 simultaneous buyers for 50 units, five runs per database. The counts were the same in
    all five runs; the times are the median of the five. One laptop, app and database on the same machine.}
  \label{tab:m1-race}
\end{table}

The counts are the result. The times are not what a customer would see in normal use: all 200 requests arrive in
the same instant and the server works on 40 at a time, so most of each time is spent waiting in a queue. The laptop
also ran the database library without its compiled speed-ups. A single hold on an idle server takes a few
milliseconds.

\subsection{What PostgreSQL found that SQLite hid}

Moving to the HLD's database was worth doing for the bugs it exposed. SQLite allows one writer at a time, which hides
mistakes that a database with real concurrency does not forgive.

\begin{itemize}
  \item \textbf{A lost update on bin quantities.} The totals row was always changed by a conditional
    \texttt{UPDATE}, but a bin's quantity was read, changed in the program and written back. With two packers
    taking from the same bin, both read the old number and one sale's change was lost: a simulated run of 100 orders
    ended with 27 units in the totals and 28 in the bins. The audit caught it. Every bin quantity is now changed
    inside the database, with the same kind of conditional \texttt{UPDATE}, and a test reproduces the case on both
    databases.
  \item \textbf{A stall under a burst.} With 200 buyers at once, requests waited 30 to 60 seconds and many failed,
    although nothing was oversold. The app kept at most 15 database connections for 40 request threads. Threads
    waited for a connection while the finished requests that would return one waited for a thread. The pool is now
    never smaller than the number of threads, and the reason is written beside the setting.
  \item \textbf{A reset that could wait forever.} PostgreSQL will not drop a table while another session still has
    a transaction open on it. Statements now give up after 15 seconds instead of waiting without end, and the
    reset of the demo data drops and recreates the tables in one transaction.
\end{itemize}

\subsection{Reads: the ATP cache}

HLD~3 expects stock reads to outnumber writes by about 13 to 1, and HLD~9.2 serves ATP from a cache. The prototype
keeps each store's availability list in the app's memory for 2 seconds and drops it whenever that store's stock
changes. Table~\ref{tab:m1-cache} shows what that saves.

\begin{table}[ht]
  \centering
  \footnotesize
  \begin{tabular}{@{}llrrr@{}}
    \toprule
    Database & Cache & Median (ms) & p99 (ms) & Hit rate \\
    \midrule
    SQLite & off & 7.64 & 10.45 & 0\% \\
    SQLite & on (2 s) & 3.93 & 8.78 & 92\% \\
    \addlinespace[3pt]
    PostgreSQL 16 & off & 12.26 & 15.85 & 0\% \\
    PostgreSQL 16 & on (2 s) & 3.85 & 13.74 & 92\% \\
    \bottomrule
  \end{tabular}
  \caption{2,000 whole-store availability reads with a checkout hold after every 13th read
    (\texttt{scripts/bench\_availability.py}). Measured inside the app's process, so the times show the database work
    saved, not what a phone would see.}
  \label{tab:m1-cache}
\end{table}

The median read is about half as long on SQLite and about a third as long on PostgreSQL, where every miss is a
round trip to another process. The p99 hardly moves, because the slowest reads are the misses that follow a
write.

\paragraph{Why a stale cache is safe.} The cache is only ever used to \emph{show} stock. The hold never reads it;
it always runs the conditional \texttt{UPDATE} against the database. A cached list can be out of date in one
window: a request that read the old numbers just before a write commits can store them just after the cache was
cleared. The 2-second expiry bounds that. In that window a customer may see an item that then fails to hold
(``sorry, just sold out''), which is an inconvenience. The store can never promise a unit it does not have, because
the promise is not made from the cache.

\subsection{Scaling out}

Sharding by store (HLD~7) is not built. It stays possible because every stock query already filters by
\texttt{store\_id} and no hold ever touches two stores, so each store's rows could move to their own database
without changing a transaction. Two things would have to change first for several app instances: the cache is per
process, so another instance can be up to 2 seconds behind, which a shared Redis would fix; and the timed jobs
should run in one instance only.

\section{Trade-off: writing off expired stock that a paid order still holds (HLD 13)}
\label{sec:m1-tradeoff}

The hourly expiry job can find an expired batch whose units are held for an order that is paid but not yet picked.
There are two choices. The job could cancel or change that order itself, which keeps the numbers tidy but makes
Inventory reach into the order flow that Fulfilment owns, from a timed job with nobody watching. Or it could write
the units off and leave the hold alone, which is what the prototype does.

The cost is a short period in which more units are held than exist for that product. The audit reports it as
\texttt{reserved\_over\_on\_hand}, and it is the one mismatch that can be legitimate. It ends when the picker reaches
the line: allocation finds no unexpired units, and the existing empty-shelf path drops the line, frees the hold and
publishes the refund event. One path handles every kind of missing item, and the customer is told by the service
that owns the order. We accepted a number that is briefly untidy in exchange for one owner per decision.

Table~\ref{tab:m1-limits} lists what the inventory work leaves open.

\begin{table}[ht]
  \centering
  \footnotesize
  \begin{tabular}{@{}>{\raggedright\arraybackslash}p{0.27\linewidth}>{\raggedright\arraybackslash}p{0.36\linewidth}
                  >{\raggedright\arraybackslash}p{0.31\linewidth}@{}}
    \toprule
    Limit & What it means & Next step \\
    \midrule
    No sharding & All stores share one database & Split by \texttt{store\_id}; the queries already allow it \\
    \addlinespace[3pt]
    Cache is per app process & A second instance can show stock up to 2 seconds old & A shared Redis with the same
      four operations \\
    \addlinespace[3pt]
    Goods receipts take no idempotency key & A receipt sent twice is received twice & Wrap the endpoint as stock
      adjustments are \\
    \addlinespace[3pt]
    Saved idempotency keys are never removed & The table grows with every keyed write & Delete keys older than a
      day \\
    \addlinespace[3pt]
    \texttt{swap\_line} and \texttt{move} have no endpoint yet & A substitute cannot be accepted and stock cannot
      be re-slotted from a screen & The fulfilment and slotting endpoints that call them \\
    \addlinespace[3pt]
    Race measured on one laptop & The times say little about a cluster & Repeat against the hosted deployment \\
    \bottomrule
  \end{tabular}
  \caption{Limits of the inventory work and what would close them.}
  \label{tab:m1-limits}
\end{table}
. Needs: array, booktabs.
% member1-standalone.tex compiles these sections on their own.
% Numbers measured on branch m1/postgres-audit-load-test (on top of main 7c95078), 21 November 2026, one laptop.

\section{Inventory (HLD 5, service 3)}
\label{sec:m1-inventory}

Every other service trusts one number from Inventory: how many units of a product a store can still sell. If that
number is too high, a customer pays for something that is not there. Inventory therefore answers two questions and
guards one rule.

\begin{itemize}
  \item \textbf{How many can we sell?} Available to promise (ATP) is on hand, minus held, minus a safety buffer, and
    never below zero. The buffer hides a few units of a product from customers so that a small counting error does
    not become a failed order. Replenishment and picking ignore it, because those units are really on the shelf.
  \item \textbf{Where are they?} Each batch of a product in each bin is one row with its expiry date. Picking takes the
    batch that expires first, then the nearest bin, and never a batch that has already expired.
  \item \textbf{The rule: a unit is promised to one customer only.} Checkout \emph{holds} units, packing
    \emph{commits} the sale, and everything in between only has to give a hold back.
\end{itemize}

\paragraph{Hold, then commit.} Checkout raises \texttt{qty\_reserved} and leaves \texttt{qty\_on\_hand} alone. Only
packing lowers both. A failed payment, a cancelled order or an expired hold therefore never has to ``put stock
back''; it lowers \texttt{qty\_reserved} and nothing else. A hold lasts 10 minutes. A timed job releases expired
holds every 30 seconds, and a hold is released by one conditional \texttt{UPDATE} on its status, so a cancel, a
payment and the timed job arriving together cannot give the same units back twice.

\paragraph{What else the service does.}
\begin{itemize}
  \item \textbf{Expiry.} An hourly job writes off every batch past its date, with one ledger row per batch, and
    publishes \texttt{item.near\_expiry} for batches within two days of their date, once per batch per day.
  \item \textbf{Idempotency keys.} A stock adjustment sent with an \texttt{Idempotency-Key} header changes stock
    once; a retry gets the first answer back. The key is saved in the same transaction as the stock change, so a
    crash between the two cannot separate them.
  \item \textbf{Cycle counts.} A counter is told which bin to count but not what the records expect. A difference
    worth at most Rs~500 is applied at once; a larger one waits for a manager, who may not approve their own count.
    When a picker reports an empty shelf, a count of that bin is opened automatically. The share of counted lines that
    matched is the inventory accuracy figure of HLD~12.
  \item \textbf{Functions for the other services.} Fulfilment asks Inventory to choose bins, to commit a sale and to
    write off units a picker could not find. Inbound asks it to receive stock. Two more were added for work the
    other members have planned: \texttt{swap\_line} holds a substitute for a dropped line, and \texttt{move}
    shifts units between two bins of the same temperature zone. No other module writes to the stock tables.
\end{itemize}

\section{Data design (HLD 7)}
\label{sec:m1-data}

\paragraph{Three records that must agree.} Stock is written in three places, and
Table~\ref{tab:m1-stock-records} says what each is for. For every store and product, the total in
\texttt{stock\_levels}, the sum of the bins in \texttt{stock\_items} and the sum of the movements in
\texttt{stock\_ledger} must be equal. Every function that changes one changes all three in the same transaction.

\begin{table}[ht]
  \centering
  \footnotesize
  \begin{tabular}{@{}>{\raggedright\arraybackslash}p{0.17\linewidth}>{\raggedright\arraybackslash}p{0.25\linewidth}
                  >{\raggedright\arraybackslash}p{0.5\linewidth}@{}}
    \toprule
    Table & One row per & Used for \\
    \midrule
    \texttt{stock\_levels} & store and product & On hand, held and safety buffer. Every checkout hold is one
      conditional \texttt{UPDATE} on this row \\
    \addlinespace[3pt]
    \texttt{stock\_items} & store, product, bin and batch & Where the units are and when they expire: pick paths,
      first-expiry-first-out, putaway, cycle counts \\
    \addlinespace[3pt]
    \texttt{stock\_ledger} & stock movement & Append-only history with a reason (receipt, sale, damage, expiry,
      theft, count, not found, move). A mistake is undone by a new row with the opposite sign, never by an edit \\
    \bottomrule
  \end{tabular}
  \caption{The three records of stock.}
  \label{tab:m1-stock-records}
\end{table}

HLD~7.1 keeps the held quantity on every batch row. The prototype keeps it on a separate totals row instead. A
checkout does not care which bin a unit is in, so the hold should lock one small row per product, not several batch
rows in an order that could deadlock. The cost is that the totals row repeats what the batch rows already say, and
two copies can drift apart.

\paragraph{The audit.} \texttt{GET /v1/admin/stock-audit} compares the three records for every store and product
and lists each disagreement with its values: level against bins, level against ledger, a negative quantity, or
more units held than exist. An empty list means the stock is consistent. The automated tests call the same function
after every test that moves stock, so a change that breaks the rule fails in review, not in a store.

\paragraph{PostgreSQL and migrations.} HLD~7.2 chooses PostgreSQL because its transactions and row locks prevent
overselling. The prototype now runs on either database, chosen by one setting, \texttt{DARKSTORE\_DB}: a SQLite file
for a laptop with nothing to install, PostgreSQL for a hosted deployment. All 313 automated tests pass on both
(SQLite and PostgreSQL~16), and the build server runs them on both for every change. Table changes are Alembic
migrations: a change to the models is turned into a migration file with one command and applied with a second, and
the app applies pending migrations itself when it starts, so a hosted database is changed in place and keeps its
data. A test fails when the models and the migrations disagree.

\section{Scalability and consistency (HLD 9)}
\label{sec:m1-consistency}

\subsection{No overselling: check and change in one statement}

The hold is a single SQL statement whose \texttt{WHERE} clause is the stock check:
\begin{center}\small
  \verb|UPDATE stock_levels SET qty_reserved = qty_reserved + :qty|\\
  \verb|WHERE store_id = :store AND sku_id = :sku|\\
  \verb|  AND qty_on_hand - qty_reserved - safety_buffer >= :qty|
\end{center}
If the statement changes no row, there was not enough stock and the whole cart is rolled back, so a cart is held
whole or not at all. Reading the level first and updating it afterwards would be wrong: two requests could both read
``1 left'' and both take it.

The statement is safe on PostgreSQL at its default isolation level (read committed) for a specific reason. When two
requests update the same row, the second waits for the first to commit, and then \emph{checks its \texttt{WHERE}
clause again against the new row} before changing anything. The second buyer of the last unit therefore sees zero
available and changes no row. SQLite reaches the same result more bluntly: it allows one writer at a time.

\subsection{The proof: 200 buyers, 50 units}

\texttt{scripts/load\_checkout.py} resets the data, puts exactly 50 units of one product on the shelf, and lets 200
customers each try to hold one unit at the same instant against a running server. It passes only if exactly 50 holds
succeed, the other 150 buyers are told \texttt{409 insufficient\_stock}, and the audit afterwards is empty.
Table~\ref{tab:m1-race} gives the result of five runs on each database.

\begin{table}[ht]
  \centering
  \footnotesize
  \begin{tabular}{@{}lrrrrrr@{}}
    \toprule
    Database & Holds & Refused & Errors & Audit mismatches & p50 (ms) & p99 (ms) \\
    \midrule
    SQLite (one file, WAL mode) & 50 & 150 & 0 & 0 & 883 & 1998 \\
    PostgreSQL 16 & 50 & 150 & 0 & 0 & 1337 & 1478 \\
    \bottomrule
  \end{tabular}
  \caption{Checkout race: 200 simultaneous buyers for 50 units, five runs per database. The counts were the same in
    all five runs; the times are the median of the five. One laptop, app and database on the same machine.}
  \label{tab:m1-race}
\end{table}

The counts are the result. The times are not what a customer would see in normal use: all 200 requests arrive in
the same instant and the server works on 40 at a time, so most of each time is spent waiting in a queue. The laptop
also ran the database library without its compiled speed-ups. A single hold on an idle server takes a few
milliseconds.

\subsection{What PostgreSQL found that SQLite hid}

Moving to the HLD's database was worth doing for the bugs it exposed. SQLite allows one writer at a time, which hides
mistakes that a database with real concurrency does not forgive.

\begin{itemize}
  \item \textbf{A lost update on bin quantities.} The totals row was always changed by a conditional
    \texttt{UPDATE}, but a bin's quantity was read, changed in the program and written back. With two packers
    taking from the same bin, both read the old number and one sale's change was lost: a simulated run of 100 orders
    ended with 27 units in the totals and 28 in the bins. The audit caught it. Every bin quantity is now changed
    inside the database, with the same kind of conditional \texttt{UPDATE}, and a test reproduces the case on both
    databases.
  \item \textbf{A stall under a burst.} With 200 buyers at once, requests waited 30 to 60 seconds and many failed,
    although nothing was oversold. The app kept at most 15 database connections for 40 request threads. Threads
    waited for a connection while the finished requests that would return one waited for a thread. The pool is now
    never smaller than the number of threads, and the reason is written beside the setting.
  \item \textbf{A reset that could wait forever.} PostgreSQL will not drop a table while another session still has
    a transaction open on it. Statements now give up after 15 seconds instead of waiting without end, and the
    reset of the demo data drops and recreates the tables in one transaction.
\end{itemize}

\subsection{Reads: the ATP cache}

HLD~3 expects stock reads to outnumber writes by about 13 to 1, and HLD~9.2 serves ATP from a cache. The prototype
keeps each store's availability list in the app's memory for 2 seconds and drops it whenever that store's stock
changes. Table~\ref{tab:m1-cache} shows what that saves.

\begin{table}[ht]
  \centering
  \footnotesize
  \begin{tabular}{@{}llrrr@{}}
    \toprule
    Database & Cache & Median (ms) & p99 (ms) & Hit rate \\
    \midrule
    SQLite & off & 7.64 & 10.45 & 0\% \\
    SQLite & on (2 s) & 3.93 & 8.78 & 92\% \\
    \addlinespace[3pt]
    PostgreSQL 16 & off & 12.26 & 15.85 & 0\% \\
    PostgreSQL 16 & on (2 s) & 3.85 & 13.74 & 92\% \\
    \bottomrule
  \end{tabular}
  \caption{2,000 whole-store availability reads with a checkout hold after every 13th read
    (\texttt{scripts/bench\_availability.py}). Measured inside the app's process, so the times show the database work
    saved, not what a phone would see.}
  \label{tab:m1-cache}
\end{table}

The median read is about half as long on SQLite and about a third as long on PostgreSQL, where every miss is a
round trip to another process. The p99 hardly moves, because the slowest reads are the misses that follow a
write.

\paragraph{Why a stale cache is safe.} The cache is only ever used to \emph{show} stock. The hold never reads it;
it always runs the conditional \texttt{UPDATE} against the database. A cached list can be out of date in one
window: a request that read the old numbers just before a write commits can store them just after the cache was
cleared. The 2-second expiry bounds that. In that window a customer may see an item that then fails to hold
(``sorry, just sold out''), which is an inconvenience. The store can never promise a unit it does not have, because
the promise is not made from the cache.

\subsection{Scaling out}

Sharding by store (HLD~7) is not built. It stays possible because every stock query already filters by
\texttt{store\_id} and no hold ever touches two stores, so each store's rows could move to their own database
without changing a transaction. Two things would have to change first for several app instances: the cache is per
process, so another instance can be up to 2 seconds behind, which a shared Redis would fix; and the timed jobs
should run in one instance only.

\section{Trade-off: writing off expired stock that a paid order still holds (HLD 13)}
\label{sec:m1-tradeoff}

The hourly expiry job can find an expired batch whose units are held for an order that is paid but not yet picked.
There are two choices. The job could cancel or change that order itself, which keeps the numbers tidy but makes
Inventory reach into the order flow that Fulfilment owns, from a timed job with nobody watching. Or it could write
the units off and leave the hold alone, which is what the prototype does.

The cost is a short period in which more units are held than exist for that product. The audit reports it as
\texttt{reserved\_over\_on\_hand}, and it is the one mismatch that can be legitimate. It ends when the picker reaches
the line: allocation finds no unexpired units, and the existing empty-shelf path drops the line, frees the hold and
publishes the refund event. One path handles every kind of missing item, and the customer is told by the service
that owns the order. We accepted a number that is briefly untidy in exchange for one owner per decision.

Table~\ref{tab:m1-limits} lists what the inventory work leaves open.

\begin{table}[ht]
  \centering
  \footnotesize
  \begin{tabular}{@{}>{\raggedright\arraybackslash}p{0.27\linewidth}>{\raggedright\arraybackslash}p{0.36\linewidth}
                  >{\raggedright\arraybackslash}p{0.31\linewidth}@{}}
    \toprule
    Limit & What it means & Next step \\
    \midrule
    No sharding & All stores share one database & Split by \texttt{store\_id}; the queries already allow it \\
    \addlinespace[3pt]
    Cache is per app process & A second instance can show stock up to 2 seconds old & A shared Redis with the same
      four operations \\
    \addlinespace[3pt]
    Goods receipts take no idempotency key & A receipt sent twice is received twice & Wrap the endpoint as stock
      adjustments are \\
    \addlinespace[3pt]
    Saved idempotency keys are never removed & The table grows with every keyed write & Delete keys older than a
      day \\
    \addlinespace[3pt]
    \texttt{swap\_line} and \texttt{move} have no endpoint yet & A substitute cannot be accepted and stock cannot
      be re-slotted from a screen & The fulfilment and slotting endpoints that call them \\
    \addlinespace[3pt]
    Race measured on one laptop & The times say little about a cluster & Repeat against the hosted deployment \\
    \bottomrule
  \end{tabular}
  \caption{Limits of the inventory work and what would close them.}
  \label{tab:m1-limits}
\end{table}
. Needs: array, booktabs.
% member1-standalone.tex compiles these sections on their own.
% Numbers measured on branch m1/postgres-audit-load-test (on top of main 7c95078), 21 November 2026, one laptop.

\section{Inventory (HLD 5, service 3)}
\label{sec:m1-inventory}

Every other service trusts one number from Inventory: how many units of a product a store can still sell. If that
number is too high, a customer pays for something that is not there. Inventory therefore answers two questions and
guards one rule.

\begin{itemize}
  \item \textbf{How many can we sell?} Available to promise (ATP) is on hand, minus held, minus a safety buffer, and
    never below zero. The buffer hides a few units of a product from customers so that a small counting error does
    not become a failed order. Replenishment and picking ignore it, because those units are really on the shelf.
  \item \textbf{Where are they?} Each batch of a product in each bin is one row with its expiry date. Picking takes the
    batch that expires first, then the nearest bin, and never a batch that has already expired.
  \item \textbf{The rule: a unit is promised to one customer only.} Checkout \emph{holds} units, packing
    \emph{commits} the sale, and everything in between only has to give a hold back.
\end{itemize}

\paragraph{Hold, then commit.} Checkout raises \texttt{qty\_reserved} and leaves \texttt{qty\_on\_hand} alone. Only
packing lowers both. A failed payment, a cancelled order or an expired hold therefore never has to ``put stock
back''; it lowers \texttt{qty\_reserved} and nothing else. A hold lasts 10 minutes. A timed job releases expired
holds every 30 seconds, and a hold is released by one conditional \texttt{UPDATE} on its status, so a cancel, a
payment and the timed job arriving together cannot give the same units back twice.

\paragraph{What else the service does.}
\begin{itemize}
  \item \textbf{Expiry.} An hourly job writes off every batch past its date, with one ledger row per batch, and
    publishes \texttt{item.near\_expiry} for batches within two days of their date, once per batch per day.
  \item \textbf{Idempotency keys.} A stock adjustment sent with an \texttt{Idempotency-Key} header changes stock
    once; a retry gets the first answer back. The key is saved in the same transaction as the stock change, so a
    crash between the two cannot separate them.
  \item \textbf{Cycle counts.} A counter is told which bin to count but not what the records expect. A difference
    worth at most Rs~500 is applied at once; a larger one waits for a manager, who may not approve their own count.
    When a picker reports an empty shelf, a count of that bin is opened automatically. The share of counted lines that
    matched is the inventory accuracy figure of HLD~12.
  \item \textbf{Functions for the other services.} Fulfilment asks Inventory to choose bins, to commit a sale and to
    write off units a picker could not find. Inbound asks it to receive stock. Two more were added for work the
    other members have planned: \texttt{swap\_line} holds a substitute for a dropped line, and \texttt{move}
    shifts units between two bins of the same temperature zone. No other module writes to the stock tables.
\end{itemize}

\section{Data design (HLD 7)}
\label{sec:m1-data}

\paragraph{Three records that must agree.} Stock is written in three places, and
Table~\ref{tab:m1-stock-records} says what each is for. For every store and product, the total in
\texttt{stock\_levels}, the sum of the bins in \texttt{stock\_items} and the sum of the movements in
\texttt{stock\_ledger} must be equal. Every function that changes one changes all three in the same transaction.

\begin{table}[ht]
  \centering
  \footnotesize
  \begin{tabular}{@{}>{\raggedright\arraybackslash}p{0.17\linewidth}>{\raggedright\arraybackslash}p{0.25\linewidth}
                  >{\raggedright\arraybackslash}p{0.5\linewidth}@{}}
    \toprule
    Table & One row per & Used for \\
    \midrule
    \texttt{stock\_levels} & store and product & On hand, held and safety buffer. Every checkout hold is one
      conditional \texttt{UPDATE} on this row \\
    \addlinespace[3pt]
    \texttt{stock\_items} & store, product, bin and batch & Where the units are and when they expire: pick paths,
      first-expiry-first-out, putaway, cycle counts \\
    \addlinespace[3pt]
    \texttt{stock\_ledger} & stock movement & Append-only history with a reason (receipt, sale, damage, expiry,
      theft, count, not found, move). A mistake is undone by a new row with the opposite sign, never by an edit \\
    \bottomrule
  \end{tabular}
  \caption{The three records of stock.}
  \label{tab:m1-stock-records}
\end{table}

HLD~7.1 keeps the held quantity on every batch row. The prototype keeps it on a separate totals row instead. A
checkout does not care which bin a unit is in, so the hold should lock one small row per product, not several batch
rows in an order that could deadlock. The cost is that the totals row repeats what the batch rows already say, and
two copies can drift apart.

\paragraph{The audit.} \texttt{GET /v1/admin/stock-audit} compares the three records for every store and product
and lists each disagreement with its values: level against bins, level against ledger, a negative quantity, or
more units held than exist. An empty list means the stock is consistent. The automated tests call the same function
after every test that moves stock, so a change that breaks the rule fails in review, not in a store.

\paragraph{PostgreSQL and migrations.} HLD~7.2 chooses PostgreSQL because its transactions and row locks prevent
overselling. The prototype now runs on either database, chosen by one setting, \texttt{DARKSTORE\_DB}: a SQLite file
for a laptop with nothing to install, PostgreSQL for a hosted deployment. All 313 automated tests pass on both
(SQLite and PostgreSQL~16), and the build server runs them on both for every change. Table changes are Alembic
migrations: a change to the models is turned into a migration file with one command and applied with a second, and
the app applies pending migrations itself when it starts, so a hosted database is changed in place and keeps its
data. A test fails when the models and the migrations disagree.

\section{Scalability and consistency (HLD 9)}
\label{sec:m1-consistency}

\subsection{No overselling: check and change in one statement}

The hold is a single SQL statement whose \texttt{WHERE} clause is the stock check:
\begin{center}\small
  \verb|UPDATE stock_levels SET qty_reserved = qty_reserved + :qty|\\
  \verb|WHERE store_id = :store AND sku_id = :sku|\\
  \verb|  AND qty_on_hand - qty_reserved - safety_buffer >= :qty|
\end{center}
If the statement changes no row, there was not enough stock and the whole cart is rolled back, so a cart is held
whole or not at all. Reading the level first and updating it afterwards would be wrong: two requests could both read
``1 left'' and both take it.

The statement is safe on PostgreSQL at its default isolation level (read committed) for a specific reason. When two
requests update the same row, the second waits for the first to commit, and then \emph{checks its \texttt{WHERE}
clause again against the new row} before changing anything. The second buyer of the last unit therefore sees zero
available and changes no row. SQLite reaches the same result more bluntly: it allows one writer at a time.

\subsection{The proof: 200 buyers, 50 units}

\texttt{scripts/load\_checkout.py} resets the data, puts exactly 50 units of one product on the shelf, and lets 200
customers each try to hold one unit at the same instant against a running server. It passes only if exactly 50 holds
succeed, the other 150 buyers are told \texttt{409 insufficient\_stock}, and the audit afterwards is empty.
Table~\ref{tab:m1-race} gives the result of five runs on each database.

\begin{table}[ht]
  \centering
  \footnotesize
  \begin{tabular}{@{}lrrrrrr@{}}
    \toprule
    Database & Holds & Refused & Errors & Audit mismatches & p50 (ms) & p99 (ms) \\
    \midrule
    SQLite (one file, WAL mode) & 50 & 150 & 0 & 0 & 883 & 1998 \\
    PostgreSQL 16 & 50 & 150 & 0 & 0 & 1337 & 1478 \\
    \bottomrule
  \end{tabular}
  \caption{Checkout race: 200 simultaneous buyers for 50 units, five runs per database. The counts were the same in
    all five runs; the times are the median of the five. One laptop, app and database on the same machine.}
  \label{tab:m1-race}
\end{table}

The counts are the result. The times are not what a customer would see in normal use: all 200 requests arrive in
the same instant and the server works on 40 at a time, so most of each time is spent waiting in a queue. The laptop
also ran the database library without its compiled speed-ups. A single hold on an idle server takes a few
milliseconds.

\subsection{What PostgreSQL found that SQLite hid}

Moving to the HLD's database was worth doing for the bugs it exposed. SQLite allows one writer at a time, which hides
mistakes that a database with real concurrency does not forgive.

\begin{itemize}
  \item \textbf{A lost update on bin quantities.} The totals row was always changed by a conditional
    \texttt{UPDATE}, but a bin's quantity was read, changed in the program and written back. With two packers
    taking from the same bin, both read the old number and one sale's change was lost: a simulated run of 100 orders
    ended with 27 units in the totals and 28 in the bins. The audit caught it. Every bin quantity is now changed
    inside the database, with the same kind of conditional \texttt{UPDATE}, and a test reproduces the case on both
    databases.
  \item \textbf{A stall under a burst.} With 200 buyers at once, requests waited 30 to 60 seconds and many failed,
    although nothing was oversold. The app kept at most 15 database connections for 40 request threads. Threads
    waited for a connection while the finished requests that would return one waited for a thread. The pool is now
    never smaller than the number of threads, and the reason is written beside the setting.
  \item \textbf{A reset that could wait forever.} PostgreSQL will not drop a table while another session still has
    a transaction open on it. Statements now give up after 15 seconds instead of waiting without end, and the
    reset of the demo data drops and recreates the tables in one transaction.
\end{itemize}

\subsection{Reads: the ATP cache}

HLD~3 expects stock reads to outnumber writes by about 13 to 1, and HLD~9.2 serves ATP from a cache. The prototype
keeps each store's availability list in the app's memory for 2 seconds and drops it whenever that store's stock
changes. Table~\ref{tab:m1-cache} shows what that saves.

\begin{table}[ht]
  \centering
  \footnotesize
  \begin{tabular}{@{}llrrr@{}}
    \toprule
    Database & Cache & Median (ms) & p99 (ms) & Hit rate \\
    \midrule
    SQLite & off & 7.64 & 10.45 & 0\% \\
    SQLite & on (2 s) & 3.93 & 8.78 & 92\% \\
    \addlinespace[3pt]
    PostgreSQL 16 & off & 12.26 & 15.85 & 0\% \\
    PostgreSQL 16 & on (2 s) & 3.85 & 13.74 & 92\% \\
    \bottomrule
  \end{tabular}
  \caption{2,000 whole-store availability reads with a checkout hold after every 13th read
    (\texttt{scripts/bench\_availability.py}). Measured inside the app's process, so the times show the database work
    saved, not what a phone would see.}
  \label{tab:m1-cache}
\end{table}

The median read is about half as long on SQLite and about a third as long on PostgreSQL, where every miss is a
round trip to another process. The p99 hardly moves, because the slowest reads are the misses that follow a
write.

\paragraph{Why a stale cache is safe.} The cache is only ever used to \emph{show} stock. The hold never reads it;
it always runs the conditional \texttt{UPDATE} against the database. A cached list can be out of date in one
window: a request that read the old numbers just before a write commits can store them just after the cache was
cleared. The 2-second expiry bounds that. In that window a customer may see an item that then fails to hold
(``sorry, just sold out''), which is an inconvenience. The store can never promise a unit it does not have, because
the promise is not made from the cache.

\subsection{Scaling out}

Sharding by store (HLD~7) is not built. It stays possible because every stock query already filters by
\texttt{store\_id} and no hold ever touches two stores, so each store's rows could move to their own database
without changing a transaction. Two things would have to change first for several app instances: the cache is per
process, so another instance can be up to 2 seconds behind, which a shared Redis would fix; and the timed jobs
should run in one instance only.

\section{Trade-off: writing off expired stock that a paid order still holds (HLD 13)}
\label{sec:m1-tradeoff}

The hourly expiry job can find an expired batch whose units are held for an order that is paid but not yet picked.
There are two choices. The job could cancel or change that order itself, which keeps the numbers tidy but makes
Inventory reach into the order flow that Fulfilment owns, from a timed job with nobody watching. Or it could write
the units off and leave the hold alone, which is what the prototype does.

The cost is a short period in which more units are held than exist for that product. The audit reports it as
\texttt{reserved\_over\_on\_hand}, and it is the one mismatch that can be legitimate. It ends when the picker reaches
the line: allocation finds no unexpired units, and the existing empty-shelf path drops the line, frees the hold and
publishes the refund event. One path handles every kind of missing item, and the customer is told by the service
that owns the order. We accepted a number that is briefly untidy in exchange for one owner per decision.

Table~\ref{tab:m1-limits} lists what the inventory work leaves open.

\begin{table}[ht]
  \centering
  \footnotesize
  \begin{tabular}{@{}>{\raggedright\arraybackslash}p{0.27\linewidth}>{\raggedright\arraybackslash}p{0.36\linewidth}
                  >{\raggedright\arraybackslash}p{0.31\linewidth}@{}}
    \toprule
    Limit & What it means & Next step \\
    \midrule
    No sharding & All stores share one database & Split by \texttt{store\_id}; the queries already allow it \\
    \addlinespace[3pt]
    Cache is per app process & A second instance can show stock up to 2 seconds old & A shared Redis with the same
      four operations \\
    \addlinespace[3pt]
    Goods receipts take no idempotency key & A receipt sent twice is received twice & Wrap the endpoint as stock
      adjustments are \\
    \addlinespace[3pt]
    Saved idempotency keys are never removed & The table grows with every keyed write & Delete keys older than a
      day \\
    \addlinespace[3pt]
    \texttt{swap\_line} and \texttt{move} have no endpoint yet & A substitute cannot be accepted and stock cannot
      be re-slotted from a screen & The fulfilment and slotting endpoints that call them \\
    \addlinespace[3pt]
    Race measured on one laptop & The times say little about a cluster & Repeat against the hosted deployment \\
    \bottomrule
  \end{tabular}
  \caption{Limits of the inventory work and what would close them.}
  \label{tab:m1-limits}
\end{table}
. Needs: array, booktabs.
% member1-standalone.tex compiles these sections on their own.
% Numbers measured on branch m1/postgres-audit-load-test (on top of main 7c95078), 21 November 2026, one laptop.

\section{Inventory (HLD 5, service 3)}
\label{sec:m1-inventory}

Every other service trusts one number from Inventory: how many units of a product a store can still sell. If that
number is too high, a customer pays for something that is not there. Inventory therefore answers two questions and
guards one rule.

\begin{itemize}
  \item \textbf{How many can we sell?} Available to promise (ATP) is on hand, minus held, minus a safety buffer, and
    never below zero. The buffer hides a few units of a product from customers so that a small counting error does
    not become a failed order. Replenishment and picking ignore it, because those units are really on the shelf.
  \item \textbf{Where are they?} Each batch of a product in each bin is one row with its expiry date. Picking takes the
    batch that expires first, then the nearest bin, and never a batch that has already expired.
  \item \textbf{The rule: a unit is promised to one customer only.} Checkout \emph{holds} units, packing
    \emph{commits} the sale, and everything in between only has to give a hold back.
\end{itemize}

\paragraph{Hold, then commit.} Checkout raises \texttt{qty\_reserved} and leaves \texttt{qty\_on\_hand} alone. Only
packing lowers both. A failed payment, a cancelled order or an expired hold therefore never has to ``put stock
back''; it lowers \texttt{qty\_reserved} and nothing else. A hold lasts 10 minutes. A timed job releases expired
holds every 30 seconds, and a hold is released by one conditional \texttt{UPDATE} on its status, so a cancel, a
payment and the timed job arriving together cannot give the same units back twice.

\paragraph{What else the service does.}
\begin{itemize}
  \item \textbf{Expiry.} An hourly job writes off every batch past its date, with one ledger row per batch, and
    publishes \texttt{item.near\_expiry} for batches within two days of their date, once per batch per day.
  \item \textbf{Idempotency keys.} A stock adjustment sent with an \texttt{Idempotency-Key} header changes stock
    once; a retry gets the first answer back. The key is saved in the same transaction as the stock change, so a
    crash between the two cannot separate them.
  \item \textbf{Cycle counts.} A counter is told which bin to count but not what the records expect. A difference
    worth at most Rs~500 is applied at once; a larger one waits for a manager, who may not approve their own count.
    When a picker reports an empty shelf, a count of that bin is opened automatically. The share of counted lines that
    matched is the inventory accuracy figure of HLD~12.
  \item \textbf{Functions for the other services.} Fulfilment asks Inventory to choose bins, to commit a sale and to
    write off units a picker could not find. Inbound asks it to receive stock. Two more were added for work the
    other members have planned: \texttt{swap\_line} holds a substitute for a dropped line, and \texttt{move}
    shifts units between two bins of the same temperature zone. No other module writes to the stock tables.
\end{itemize}

\section{Data design (HLD 7)}
\label{sec:m1-data}

\paragraph{Three records that must agree.} Stock is written in three places, and
Table~\ref{tab:m1-stock-records} says what each is for. For every store and product, the total in
\texttt{stock\_levels}, the sum of the bins in \texttt{stock\_items} and the sum of the movements in
\texttt{stock\_ledger} must be equal. Every function that changes one changes all three in the same transaction.

\begin{table}[ht]
  \centering
  \footnotesize
  \begin{tabular}{@{}>{\raggedright\arraybackslash}p{0.17\linewidth}>{\raggedright\arraybackslash}p{0.25\linewidth}
                  >{\raggedright\arraybackslash}p{0.5\linewidth}@{}}
    \toprule
    Table & One row per & Used for \\
    \midrule
    \texttt{stock\_levels} & store and product & On hand, held and safety buffer. Every checkout hold is one
      conditional \texttt{UPDATE} on this row \\
    \addlinespace[3pt]
    \texttt{stock\_items} & store, product, bin and batch & Where the units are and when they expire: pick paths,
      first-expiry-first-out, putaway, cycle counts \\
    \addlinespace[3pt]
    \texttt{stock\_ledger} & stock movement & Append-only history with a reason (receipt, sale, damage, expiry,
      theft, count, not found, move). A mistake is undone by a new row with the opposite sign, never by an edit \\
    \bottomrule
  \end{tabular}
  \caption{The three records of stock.}
  \label{tab:m1-stock-records}
\end{table}

HLD~7.1 keeps the held quantity on every batch row. The prototype keeps it on a separate totals row instead. A
checkout does not care which bin a unit is in, so the hold should lock one small row per product, not several batch
rows in an order that could deadlock. The cost is that the totals row repeats what the batch rows already say, and
two copies can drift apart.

\paragraph{The audit.} \texttt{GET /v1/admin/stock-audit} compares the three records for every store and product
and lists each disagreement with its values: level against bins, level against ledger, a negative quantity, or
more units held than exist. An empty list means the stock is consistent. The automated tests call the same function
after every test that moves stock, so a change that breaks the rule fails in review, not in a store.

\paragraph{PostgreSQL and migrations.} HLD~7.2 chooses PostgreSQL because its transactions and row locks prevent
overselling. The prototype now runs on either database, chosen by one setting, \texttt{DARKSTORE\_DB}: a SQLite file
for a laptop with nothing to install, PostgreSQL for a hosted deployment. All 313 automated tests pass on both
(SQLite and PostgreSQL~16), and the build server runs them on both for every change. Table changes are Alembic
migrations: a change to the models is turned into a migration file with one command and applied with a second, and
the app applies pending migrations itself when it starts, so a hosted database is changed in place and keeps its
data. A test fails when the models and the migrations disagree.

\section{Scalability and consistency (HLD 9)}
\label{sec:m1-consistency}

\subsection{No overselling: check and change in one statement}

The hold is a single SQL statement whose \texttt{WHERE} clause is the stock check:
\begin{center}\small
  \verb|UPDATE stock_levels SET qty_reserved = qty_reserved + :qty|\\
  \verb|WHERE store_id = :store AND sku_id = :sku|\\
  \verb|  AND qty_on_hand - qty_reserved - safety_buffer >= :qty|
\end{center}
If the statement changes no row, there was not enough stock and the whole cart is rolled back, so a cart is held
whole or not at all. Reading the level first and updating it afterwards would be wrong: two requests could both read
``1 left'' and both take it.

The statement is safe on PostgreSQL at its default isolation level (read committed) for a specific reason. When two
requests update the same row, the second waits for the first to commit, and then \emph{checks its \texttt{WHERE}
clause again against the new row} before changing anything. The second buyer of the last unit therefore sees zero
available and changes no row. SQLite reaches the same result more bluntly: it allows one writer at a time.

\subsection{The proof: 200 buyers, 50 units}

\texttt{scripts/load\_checkout.py} resets the data, puts exactly 50 units of one product on the shelf, and lets 200
customers each try to hold one unit at the same instant against a running server. It passes only if exactly 50 holds
succeed, the other 150 buyers are told \texttt{409 insufficient\_stock}, and the audit afterwards is empty.
Table~\ref{tab:m1-race} gives the result of five runs on each database.

\begin{table}[ht]
  \centering
  \footnotesize
  \begin{tabular}{@{}lrrrrrr@{}}
    \toprule
    Database & Holds & Refused & Errors & Audit mismatches & p50 (ms) & p99 (ms) \\
    \midrule
    SQLite (one file, WAL mode) & 50 & 150 & 0 & 0 & 883 & 1998 \\
    PostgreSQL 16 & 50 & 150 & 0 & 0 & 1337 & 1478 \\
    \bottomrule
  \end{tabular}
  \caption{Checkout race: 200 simultaneous buyers for 50 units, five runs per database. The counts were the same in
    all five runs; the times are the median of the five. One laptop, app and database on the same machine.}
  \label{tab:m1-race}
\end{table}

The counts are the result. The times are not what a customer would see in normal use: all 200 requests arrive in
the same instant and the server works on 40 at a time, so most of each time is spent waiting in a queue. The laptop
also ran the database library without its compiled speed-ups. A single hold on an idle server takes a few
milliseconds.

\subsection{What PostgreSQL found that SQLite hid}

Moving to the HLD's database was worth doing for the bugs it exposed. SQLite allows one writer at a time, which hides
mistakes that a database with real concurrency does not forgive.

\begin{itemize}
  \item \textbf{A lost update on bin quantities.} The totals row was always changed by a conditional
    \texttt{UPDATE}, but a bin's quantity was read, changed in the program and written back. With two packers
    taking from the same bin, both read the old number and one sale's change was lost: a simulated run of 100 orders
    ended with 27 units in the totals and 28 in the bins. The audit caught it. Every bin quantity is now changed
    inside the database, with the same kind of conditional \texttt{UPDATE}, and a test reproduces the case on both
    databases.
  \item \textbf{A stall under a burst.} With 200 buyers at once, requests waited 30 to 60 seconds and many failed,
    although nothing was oversold. The app kept at most 15 database connections for 40 request threads. Threads
    waited for a connection while the finished requests that would return one waited for a thread. The pool is now
    never smaller than the number of threads, and the reason is written beside the setting.
  \item \textbf{A reset that could wait forever.} PostgreSQL will not drop a table while another session still has
    a transaction open on it. Statements now give up after 15 seconds instead of waiting without end, and the
    reset of the demo data drops and recreates the tables in one transaction.
\end{itemize}

\subsection{Reads: the ATP cache}

HLD~3 expects stock reads to outnumber writes by about 13 to 1, and HLD~9.2 serves ATP from a cache. The prototype
keeps each store's availability list in the app's memory for 2 seconds and drops it whenever that store's stock
changes. Table~\ref{tab:m1-cache} shows what that saves.

\begin{table}[ht]
  \centering
  \footnotesize
  \begin{tabular}{@{}llrrr@{}}
    \toprule
    Database & Cache & Median (ms) & p99 (ms) & Hit rate \\
    \midrule
    SQLite & off & 7.64 & 10.45 & 0\% \\
    SQLite & on (2 s) & 3.93 & 8.78 & 92\% \\
    \addlinespace[3pt]
    PostgreSQL 16 & off & 12.26 & 15.85 & 0\% \\
    PostgreSQL 16 & on (2 s) & 3.85 & 13.74 & 92\% \\
    \bottomrule
  \end{tabular}
  \caption{2,000 whole-store availability reads with a checkout hold after every 13th read
    (\texttt{scripts/bench\_availability.py}). Measured inside the app's process, so the times show the database work
    saved, not what a phone would see.}
  \label{tab:m1-cache}
\end{table}

The median read is about half as long on SQLite and about a third as long on PostgreSQL, where every miss is a
round trip to another process. The p99 hardly moves, because the slowest reads are the misses that follow a
write.

\paragraph{Why a stale cache is safe.} The cache is only ever used to \emph{show} stock. The hold never reads it;
it always runs the conditional \texttt{UPDATE} against the database. A cached list can be out of date in one
window: a request that read the old numbers just before a write commits can store them just after the cache was
cleared. The 2-second expiry bounds that. In that window a customer may see an item that then fails to hold
(``sorry, just sold out''), which is an inconvenience. The store can never promise a unit it does not have, because
the promise is not made from the cache.

\subsection{Scaling out}

Sharding by store (HLD~7) is not built. It stays possible because every stock query already filters by
\texttt{store\_id} and no hold ever touches two stores, so each store's rows could move to their own database
without changing a transaction. Two things would have to change first for several app instances: the cache is per
process, so another instance can be up to 2 seconds behind, which a shared Redis would fix; and the timed jobs
should run in one instance only.

\section{Trade-off: writing off expired stock that a paid order still holds (HLD 13)}
\label{sec:m1-tradeoff}

The hourly expiry job can find an expired batch whose units are held for an order that is paid but not yet picked.
There are two choices. The job could cancel or change that order itself, which keeps the numbers tidy but makes
Inventory reach into the order flow that Fulfilment owns, from a timed job with nobody watching. Or it could write
the units off and leave the hold alone, which is what the prototype does.

The cost is a short period in which more units are held than exist for that product. The audit reports it as
\texttt{reserved\_over\_on\_hand}, and it is the one mismatch that can be legitimate. It ends when the picker reaches
the line: allocation finds no unexpired units, and the existing empty-shelf path drops the line, frees the hold and
publishes the refund event. One path handles every kind of missing item, and the customer is told by the service
that owns the order. We accepted a number that is briefly untidy in exchange for one owner per decision.

Table~\ref{tab:m1-limits} lists what the inventory work leaves open.

\begin{table}[ht]
  \centering
  \footnotesize
  \begin{tabular}{@{}>{\raggedright\arraybackslash}p{0.27\linewidth}>{\raggedright\arraybackslash}p{0.36\linewidth}
                  >{\raggedright\arraybackslash}p{0.31\linewidth}@{}}
    \toprule
    Limit & What it means & Next step \\
    \midrule
    No sharding & All stores share one database & Split by \texttt{store\_id}; the queries already allow it \\
    \addlinespace[3pt]
    Cache is per app process & A second instance can show stock up to 2 seconds old & A shared Redis with the same
      four operations \\
    \addlinespace[3pt]
    Goods receipts take no idempotency key & A receipt sent twice is received twice & Wrap the endpoint as stock
      adjustments are \\
    \addlinespace[3pt]
    Saved idempotency keys are never removed & The table grows with every keyed write & Delete keys older than a
      day \\
    \addlinespace[3pt]
    \texttt{swap\_line} and \texttt{move} have no endpoint yet & A substitute cannot be accepted and stock cannot
      be re-slotted from a screen & The fulfilment and slotting endpoints that call them \\
    \addlinespace[3pt]
    Race measured on one laptop & The times say little about a cluster & Repeat against the hosted deployment \\
    \bottomrule
  \end{tabular}
  \caption{Limits of the inventory work and what would close them.}
  \label{tab:m1-limits}
\end{table}
